% Source coverage: physical PDF pp. 71--76 (printed pp. 48--53), from the
% Section 25.5 heading through the final paragraph immediately before Problems.
% Inventory: Eqs. (25.5.1)--(25.5.24), four unnumbered display groups, two
% source footnotes, no citations, figures, tables, or three-asterisk dividers.
% The source's operator-adjoint asterisks are modernized to daggers; complex
% conjugation of numerical matrix elements remains denoted by an asterisk.
\subsection{Massive Particle Supermultiplets}

Although the known quarks, leptons, and gauge bosons and their superpartners
may probably be treated as massless at energies where supersymmetry breaking
is negligible, this is not necessarily true for other particles, including
the extra gauge bosons of large mass required by theories that unify the
strong and electroweak interactions. Also, ever since the Wess--Zumino model,
massive particle theories have been useful test cases for studying
supersymmetry theories. It will therefore be worthwhile for us briefly to
consider the implications of unbroken supersymmetry for massive particles.

As in the previous section, we obtain the various one-particle states in a
supermultiplet by acting on any one of them with the operators
$\WeylQ_{ar}$ and $\WeylQ_{ar}^\dagger$, and all of these states have the same
four-momentum. Unlike the case of zero mass, for mass $M>0$ we may now take
this to be the four-momentum of a particle at rest, with $p^i=0$ for
$i=1,2,3$ and $p^0=M$. In this frame of reference, we have
\begin{equation}
 \sigma_\mu p^\mu=M\sigma_0
 =M\begin{pmatrix}1&0\\0&1\end{pmatrix}.
\tag{25.5.1}\label{eq:25.5.1}
\end{equation}
Thus, acting on any state $\ket{\ }$ in a supermultiplet with this
four-momentum, the anticommutation relation \eqref{eq:25.2.7} yields
\begin{equation}
 \{\WeylQ_{ar},\WeylQ_{bs}^\dagger\}\ket{\ }
 =2M\delta_{ab}\delta_{rs}\ket{\ }.
\tag{25.5.2}\label{eq:25.5.2}
\end{equation}
In contrast with the case of zero mass, here no component of $\WeylQ_{ar}$ or
$\WeylQ_{ar}^\dagger$ can vanish on the whole multiplet, so we have two sets
of raising and lowering operators: both $\WeylQ_{(1/2)r}$ and
$\WeylQ_{(-1/2)r}^\dagger$ lower the spin 3-component by $1/2$, while both
$\WeylQ_{(-1/2)r}$ and $\WeylQ_{(1/2)r}^\dagger$ raise the spin 3-component
by $1/2$. However, as we shall see, for extended supersymmetry it is possible
for certain linear combinations of the $\WeylQ$s and $\WeylQ^\dagger$s to
vanish.

We will first consider the case of simple supersymmetry. By using the
supersymmetry algebra \eqref{eq:25.2.31}, \eqref{eq:25.2.32}, we shall show
that the general massive supermultiplet consists of a particle of spin
$j+1/2$, a pair of particles of spin $j$, and a particle of spin $j-1/2$.
Where parity is conserved, the particles of spin $j\pm1/2$ have equal
intrinsic parity, given by some phase $\eta$, while the two particles of spin
$j$ have parities $+\ii\eta$ and $-\ii\eta$. Here $j$ is any integer or
half-integer greater than zero. There is also a collapsed supermultiplet,
consisting of two particles of spin zero and a particle of spin $1/2$. When
parity is conserved the particles of spin zero have parities $\ii\eta$ and
$-\ii\eta$, where $\eta$ is the parity of the particle of spin $1/2$.

Here is the proof. We first show that any supermultiplet will contain at least
one spin multiplet of states $\ket{j,\sigma}$ with spin 3-component $\sigma$
running by unit steps from $-j$ to $+j$, having the special property that, for
all such $\sigma$ and for $a=\pm1/2$,
\begin{equation}
 \WeylQ_a\ket{j,\sigma}=0.
\tag{25.5.3}\label{eq:25.5.3}
\end{equation}
Starting with any non-zero state $\ket{\psi}$ in the supermultiplet, we can
define non-zero states
\[
 \ket{\psi'}\equiv
 \begin{cases}
 (2M)^{-1/2}\WeylQ_{1/2}\ket{\psi},
   &\WeylQ_{1/2}\ket{\psi}\ne0,\\
 \ket{\psi},
   &\WeylQ_{1/2}\ket{\psi}=0,
 \end{cases}
\]
and
\[
 \ket{\psi''}\equiv
 \begin{cases}
 (2M)^{-1/2}\WeylQ_{-1/2}\ket{\psi'},
   &\WeylQ_{-1/2}\ket{\psi'}\ne0,\\
 \ket{\psi'},
   &\WeylQ_{-1/2}\ket{\psi'}=0.
 \end{cases}
\]
Because the $\WeylQ_a$ anticommute, $\WeylQ_{1/2}\ket{\psi'}=0$, and so
$\WeylQ_a\ket{\psi''}=0$ for $a=\pm1/2$. If any state $\ket{\psi''}$
satisfies the condition that $\WeylQ_a\ket{\psi''}=0$, then so does
$U(R)\ket{\psi''}$, where $U(R)$ is the unitary operator representing an
arbitrary spatial rotation. It follows that the states that satisfy this
condition may be decomposed into complete spin multiplets $\ket{j,\sigma}$,
satisfying condition \eqref{eq:25.5.3}.

Now focus on any one of these spin multiplets satisfying
Eq.~\eqref{eq:25.5.3}, normalized so that
\begin{equation}
 \braket{j,\sigma'}{j,\sigma}=\delta_{\sigma'\sigma}.
\tag{25.5.4}\label{eq:25.5.4}
\end{equation}
For $j>0$, by applying the spin $1/2$ operators\footnote{Since $\WeylQ_a$
transforms under rotations like a field that destroys a particle of spin
$1/2$ and spin 3-component $a$, it is $\WeylQ_a^\dagger$ that transforms
like a field that creates such a particle, and hence transforms like the
particle itself. To put this formally,
$[J_i,\WeylQ_a]=-\sum_b\frac12(\sigma_i)_{ab}\WeylQ_b$, so
$[J_i,\WeylQ_a^\dagger]
=\sum_b\frac12(\sigma_i)_{ba}\WeylQ_b^\dagger$, which may be compared with
the transformation property of a spin $1/2$ particle,
$J_i\ket{a}=\sum_b\frac12(\sigma_i)_{ba}\ket{b}$.}
$\WeylQ_a^\dagger$ to these states we can construct states of spin
$j\pm1/2$:
\begin{equation}
 \ket{j\pm1/2,\sigma}
 =\frac{1}{\sqrt{2M}}\sum_a
 C_{\frac12 j}(j\pm1/2,\sigma;a,\sigma-a)
 \WeylQ_a^\dagger\ket{j,\sigma-a},
\tag{25.5.5}\label{eq:25.5.5}
\end{equation}
where $C_{jj'}(j'',\sigma'';\sigma,\sigma')$ is the conventional
Clebsch--Gordan coefficient for coupling spins $j$ and $j'$ with 3-components
$\sigma$ and $\sigma'$ to make spin $j''$ with 3-component $\sigma''$. Using
Eqs.~\eqref{eq:25.5.2}--\eqref{eq:25.5.5} and the orthonormality properties
of the Clebsch--Gordan coefficients, we can show that these states are
properly normalized:
\begin{equation}
 \braket{j\pm1/2,\sigma}{j\pm1/2,\sigma'}
 =\delta_{\sigma\sigma'},
 \qquad
 \braket{j\pm1/2,\sigma}{j\mp1/2,\sigma'}=0,
\tag{25.5.6}\label{eq:25.5.6}
\end{equation}
so none of the states $\ket{j\pm1/2,\sigma}$ can vanish. The only exception
is for $j=0$, in which case of course there is no state
$\ket{j-1/2,\sigma}$. We can also construct other states by applying
\emph{two} $\WeylQ^\dagger$s to $\ket{j,\sigma}$. Since each
$\WeylQ_a^\dagger$ anticommutes with itself, the only such non-zero states
are formed by applying the operator
$\WeylQ_{1/2}^\dagger\WeylQ_{-1/2}^\dagger
=-\WeylQ_{-1/2}^\dagger\WeylQ_{1/2}^\dagger$. This operator may be written
as $\frac12 e_{ab}\WeylQ_a^\dagger\WeylQ_b^\dagger$, which shows that it is
rotationally invariant, so this gives a second spin multiplet with spin $j$:
\begin{equation}
 \ket{j,\sigma}^{\mathrm b}
 =\frac{1}{2M}\WeylQ_{1/2}^\dagger
 \WeylQ_{-1/2}^\dagger\ket{j,\sigma},
\tag{25.5.7}\label{eq:25.5.7}
\end{equation}
which is distinguished from $\ket{j,\sigma}$ by the fact that instead of
Eq.~\eqref{eq:25.5.3}, we have
\begin{equation}
 \WeylQ_a^\dagger\ket{j,\sigma}^{\mathrm b}=0.
\tag{25.5.8}\label{eq:25.5.8}
\end{equation}
Again using Eqs.~\eqref{eq:25.5.2}--\eqref{eq:25.5.4}, we find that these are
also normalized states:
\begin{equation}
 {}^{\mathrm b}\!\braket{j,\sigma'}{j,\sigma}^{\mathrm b}
 =\delta_{\sigma'\sigma},
 \qquad
 \braket{j,\sigma'}{j,\sigma}^{\mathrm b}=0.
\tag{25.5.9}\label{eq:25.5.9}
\end{equation}
It is then easy to show that the states constructed so far form a complete
representation of the supersymmetry algebra. The orthonormality property of
the Clebsch--Gordan coefficients allows us to rewrite Eq.~\eqref{eq:25.5.5}
as
\begin{equation}
 \WeylQ_a^\dagger\ket{j,\sigma}
 =\sqrt{2M}\sum_{\pm}
 C_{\frac12 j}(j\pm1/2,\sigma+a;a,\sigma)
 \ket{j\pm1/2,\sigma+a}.
\tag{25.5.10}\label{eq:25.5.10}
\end{equation}
Also, Eq.~\eqref{eq:25.5.2} shows that, for any state $\ket{\ }$ in the
supermultiplet,
\begin{equation}
 [\WeylQ_a,\WeylQ_{1/2}^\dagger\WeylQ_{-1/2}^\dagger]\ket{\ }
 =2M\sum_b e_{ab}\WeylQ_b^\dagger\ket{\ },
\tag{25.5.11}\label{eq:25.5.11}
\end{equation}
so Eqs.~\eqref{eq:25.5.7} and \eqref{eq:25.5.3} give
\begin{equation}
\begin{aligned}
 \WeylQ_a\ket{j,\sigma}^{\mathrm b}
 &=\sum_b e_{ab}\WeylQ_b^\dagger\ket{j,\sigma}
 \\
 &=\sqrt{2M}\sum_b e_{ab}\sum_{\pm}
 C_{\frac12 j}(j\pm1/2,\sigma+b;b,\sigma)
 \ket{j\pm1/2,\sigma+b}.
\end{aligned}
\tag{25.5.12}\label{eq:25.5.12}
\end{equation}
From Eqs.~\eqref{eq:25.5.2}, \eqref{eq:25.5.3}, and \eqref{eq:25.5.5} we
have
\begin{equation}
 \WeylQ_a\ket{j\pm1/2,\sigma}
 =\sqrt{2M}\,
 C_{\frac12 j}(j\pm1/2,\sigma;a,\sigma-a)
 \ket{j,\sigma-a},
\tag{25.5.13}\label{eq:25.5.13}
\end{equation}
while Eqs.~\eqref{eq:25.5.5}, \eqref{eq:25.2.31}, and
\eqref{eq:25.5.7} yield
\begin{equation}
 \WeylQ_a^\dagger\ket{j\pm1/2,\sigma}
 =\sqrt{2M}\sum_b e_{ab}
 C_{\frac12 j}(j\pm1/2,\sigma;b,\sigma-b)
 \ket{j,\sigma-b}^{\mathrm b}.
\tag{25.5.14}\label{eq:25.5.14}
\end{equation}
Eqs.~\eqref{eq:25.5.3}, \eqref{eq:25.5.8}, \eqref{eq:25.5.10}, and
\eqref{eq:25.5.12}--\eqref{eq:25.5.14} give the action of the $\WeylQ$s and
$\WeylQ^\dagger$s on all the states of the supermultiplet.

For $j=0$ we have the collapsed supermultiplet:
Eqs.~\eqref{eq:25.5.3}, \eqref{eq:25.5.8}, \eqref{eq:25.5.10}, and
\eqref{eq:25.5.12}--\eqref{eq:25.5.14} become
\begin{equation}
\begin{alignedat}{2}
 \WeylQ_a\ket{0,0}&=0,
 &\qquad
 \WeylQ_a^\dagger\ket{0,0}^{\mathrm b}&=0,
 \\
 \WeylQ_a^\dagger\ket{0,0}&=\sqrt{2M}\ket{1/2,a},
 &
 \WeylQ_a\ket{0,0}^{\mathrm b}
 &=\sqrt{2M}\sum_b e_{ab}\ket{1/2,b},
 \\
 \WeylQ_a\ket{1/2,b}&=\sqrt{2M}\delta_{ab}\ket{0,0},
 &
 \WeylQ_a^\dagger\ket{1/2,b}
 &=\sqrt{2M}e_{ab}\ket{0,0}^{\mathrm b}.
\end{alignedat}
\tag{25.5.15}\label{eq:25.5.15}
\end{equation}

Now suppose that parity is conserved. Recall that the phase of the
supersymmetry generator may be chosen so that the action of the parity
operator on these generators is given by Eq.~\eqref{eq:25.3.13}. Then
$\WeylQ_a^\dagger$ acting on $P\ket{j,\sigma}$ is a linear combination of the
states $P\WeylQ_a\ket{j,\sigma}$, which vanish, and since
$P\ket{j,\sigma}$ has the same rotation properties as
$\ket{j,\sigma}^{\mathrm b}$, it must simply be proportional to it
\begin{equation}
 P\ket{j,\sigma}=-\eta\ket{j,\sigma}^{\mathrm b}.
\tag{25.5.16}\label{eq:25.5.16}
\end{equation}
Since $P$ is unitary, $\eta$ is a phase factor, with $|\eta|=1$. A
corresponding argument shows that $P\ket{j,\sigma}^{\mathrm b}$ is
proportional to $\ket{j,\sigma}$. To find the proportionality coefficient,
we note that
\[
\begin{aligned}
 P\ket{j,\sigma}^{\mathrm b}
 &=(2M)^{-1}P\WeylQ_{1/2}^\dagger
   \WeylQ_{-1/2}^\dagger\ket{j,\sigma}
 =-\eta(2M)^{-1}\WeylQ_{-1/2}\WeylQ_{1/2}
   \ket{j,\sigma}^{\mathrm b}
 \\
 &=-\eta(2M)^{-2}
   \WeylQ_{-1/2}\WeylQ_{1/2}
   \WeylQ_{1/2}^\dagger\WeylQ_{-1/2}^\dagger\ket{j,\sigma}
 =-\eta\ket{j,\sigma}.
\end{aligned}
\]
We can then define states of spin $j$
\begin{equation}
 \ket{j,\sigma}^{\pm}
 \equiv\frac{1}{\sqrt{2}}
 \left(\ket{j,\sigma}\pm \ii\ket{j,\sigma}^{\mathrm b}\right),
\tag{25.5.17}\label{eq:25.5.17}
\end{equation}
with definite parity
\begin{equation}
 P\ket{j,\sigma}^{\pm}=\pm \ii\eta\ket{j,\sigma}^{\pm}.
\tag{25.5.18}\label{eq:25.5.18}
\end{equation}
Finally, applying the parity operator to Eq.~\eqref{eq:25.5.5} and using
Eqs.~\eqref{eq:25.3.13} and \eqref{eq:25.5.16} gives
\[
 P\ket{j\pm1/2,\sigma}
 =-\frac{\eta}{\sqrt{2M}}\sum_a
 C_{\frac12 j}(j\pm1/2,\sigma;a,\sigma-a)
 \sum_b e_{ab}\WeylQ_b\ket{j,\sigma-a}^{\mathrm b}.
\]
Eq.~\eqref{eq:25.5.12} and the orthonormality property of the
Clebsch--Gordan coefficients then yield
\begin{equation}
 P\ket{j\pm1/2,\sigma}=\eta\ket{j\pm1/2,\sigma},
\tag{25.5.19}\label{eq:25.5.19}
\end{equation}
as was to be shown.

We now turn briefly to the case of extended supersymmetry, with $N$
supersymmetry generators. As mentioned in the previous section, there can be
no massless particle with a non-vanishing eigenvalue for any central charge.
We can go further and show that the eigenvalues of the central charge
operators set a lower bound on the mass of any supermultiplet.

Because the central charges $Z_{rs}$ and $Z_{rs}^\dagger$ commute with each
other and with $P_\mu$, one-particle states can be chosen to be eigenstates of
all the central charges as well as of $P_\mu$, and because the central charges
commute with the $\WeylQ_{ar}$ and $\WeylQ_{ar}^\dagger$, all states in a
supermultiplet have the same eigenvalues.

To derive an inequality relating the mass $M$ of a supermultiplet and the
eigenvalues of the central charges on this multiplet, we use the
anticommutation relations \eqref{eq:25.2.7} and \eqref{eq:25.2.8} to write
\begin{equation}
\begin{aligned}
 &\sum_{ar}\left\{
 \left(\WeylQ_{ar}-\sum_{bs}e_{ab}U_{rs}\WeylQ_{bs}^\dagger\right),
 \left(\WeylQ_{ar}^\dagger
       -\sum_{ct}e_{ac}U_{rt}^*\WeylQ_{ct}\right)
 \right\}
 \\
 &\hspace{4cm}
 =8NP^0-2\operatorname{Tr}\left(ZU^\dagger+UZ^\dagger\right),
\end{aligned}
\tag{25.5.20}\label{eq:25.5.20}
\end{equation}
where $U_{rs}$ is an arbitrary $N\times N$ unitary matrix. The left-hand side
is a positive-definite operator, so by letting this act on the states of the
supermultiplet at rest, we find
\begin{equation}
 M\geq\frac{1}{4N}\operatorname{Tr}
 \left(ZU^\dagger+UZ^\dagger\right),
\tag{25.5.21}\label{eq:25.5.21}
\end{equation}
where now $Z_{rs}$ denotes the values of the central charges for the
supermultiplet of mass $M$. The polar decomposition theorem tells us that any
square matrix $Z$ may be written as $HV$, where $H$ is a positive Hermitian
matrix and $V$ is unitary. We can obtain a useful inequality (which is in fact
optimal) by setting $U=V$, in which case Eq.~\eqref{eq:25.5.21} becomes
\begin{equation}
 M\geq\frac{1}{2N}\operatorname{Tr}H
 =\frac{1}{2N}\operatorname{Tr}\sqrt{Z^\dagger Z}.
\tag{25.5.22}\label{eq:25.5.22}
\end{equation}
States for which $M$ equals the minimum value allowed by this inequality are
known as BPS states, by analogy with the
Bogomol'nyi--Prasad--Sommerfeld magnetic monopole configurations discussed in
Section 23.3, whose mass also equals a lower bound on general monopole masses.
In fact, this is more than an analogy; we will see in Section 27.9 that the
lower bound on monopole masses in gauge theories with extended supersymmetry
is a special case of the lower bound \eqref{eq:25.5.22}.

As we can see from this derivation of Eq.~\eqref{eq:25.5.22}, for BPS
supermultiplets the operator
$\WeylQ_{ar}-\sum_{bs}e_{ab}U_{rs}\WeylQ_{bs}^\dagger$ gives zero when acting
on any state of a supermultiplet, so there are only $N$ independent
helicity-lowering operators $\WeylQ_{(1/2)r}$ and only $N$ independent
helicity-raising operators $\WeylQ_{(-1/2)r}$, just as in the case of massless
supermultiplets. This leads to smaller supermultiplets than would be found in
the general case.

For instance, for $N=2$ supersymmetry the central charge is given by a single
complex number\footnote{In some articles on $N=2$ supersymmetry the central
charge $Z$ is what we would call $Z/(2\sqrt{2})$.}
\begin{equation}
 Z=\begin{pmatrix}0&Z_{12}\\-Z_{12}&0\end{pmatrix}.
\tag{25.5.23}\label{eq:25.5.23}
\end{equation}
The inequality \eqref{eq:25.5.22} here reads
\begin{equation}
 M\geq|Z_{12}|/2.
\tag{25.5.24}\label{eq:25.5.24}
\end{equation}
Where $M=|Z_{12}|/2$, the helicity contents of the massive particle
supermultiplets are the same as for those of zero mass: there are gauge
supermultiplets, consisting of one particle of spin 1, an $SU(2)$
$R$-symmetry doublet of spin $1/2$, and one particle of spin 0 (the other
helicity zero state belonging to the spin one particle), and hypermultiplets,
consisting of one particle of spin $1/2$ and an $SU(2)$ $R$-symmetry doublet
of spin 0. These are sometimes called `short' supermultiplets, to distinguish
them from the larger supermultiplets encountered when $M>|Z_{12}|/2$.
